\chapter{Routing Protocols and MAC Layer Theory}
\label{ch:routing}

\section{Classification of Routing Protocols}

Routing protocols for ad~hoc and mesh networks are broadly classified into three
categories based on when and how they compute routes~\cite{abolhasan2004review}:

\subsection{Proactive (Table-Driven) Protocols}

Proactive protocols maintain a continuously updated routing table that contains entries
for all reachable destinations in the network. Whenever a topology change is detected
--- a link failure, a new node joining, or a metric update --- the protocol propagates
update messages throughout the network to synchronise all routing tables.

\textbf{Advantages:}
\begin{itemize}
  \item Routes are immediately available when a transmission session begins; no discovery delay.
  \item Predictable and bounded latency, since the routing layer never needs to wait for route discovery.
  \item Simple implementation: standard distance-vector or link-state algorithms adapted for wireless media.
\end{itemize}

\textbf{Disadvantages:}
\begin{itemize}
  \item High control overhead, especially in dense networks where many nodes exchange routing updates periodically.
  \item Routing table entries for destinations that are never communicated with consume memory and channel capacity unnecessarily.
  \item Convergence delay after topology changes can leave stale routes in the table, potentially causing black-hole or routing-loop problems.
\end{itemize}

Examples of proactive protocols include DSDV~\cite{perkins1994dsdv}, OLSR~\cite{clausen2003olsr}, and BATMAN.

\subsection{Reactive (On-Demand) Protocols}

Reactive protocols do not maintain routing tables for all destinations. Instead, they
initiate a route discovery procedure --- typically by flooding a Route Request (RREQ)
message throughout the network --- only when a node needs to communicate with a
destination for which it has no valid route. A Route Reply (RREP) is sent back from
the destination to the source along the reverse path, establishing the route. The route
is cached and used for subsequent packets to the same destination until it expires or a
link failure is detected.

\textbf{Advantages:}
\begin{itemize}
  \item Low steady-state control overhead; no periodic broadcasts when the network is idle.
  \item Scales to larger networks because routing state is built on demand.
  \item Energy-efficient: radios spend more time in the receive or sleep state rather than transmitting periodic updates.
\end{itemize}

\textbf{Disadvantages:}
\begin{itemize}
  \item Route discovery latency on first contact with a destination, which can be significant in sparse networks where the RREQ must traverse many hops.
  \item RREQ flooding can cause a broadcast storm in dense networks if not carefully controlled.
  \item Maintaining routes for long-lived sessions requires periodic HELLO messages, adding some overhead.
\end{itemize}

Examples of reactive protocols include AODV~\cite{perkins2003aodv}, DSR~\cite{johnson2007dsr}, and DYMO.

\subsection{Hybrid Protocols}

Hybrid protocols combine proactive and reactive mechanisms. They maintain routing tables
proactively within a local zone of a configurable radius, and use reactive flooding for
destinations outside the zone. The Zone Routing Protocol (ZRP) and HWMP (the protocol
mandated by IEEE~802.11s for standardised mesh networking) are prominent examples.

Hybrid protocols aim to capture the low-latency benefit of proactive routing for nearby
nodes while limiting overhead by restricting proactive maintenance to the local
neighbourhood. However, the zone radius must be tuned carefully to achieve the desired
trade-off, introducing an additional configuration parameter.

Table~\ref{tab:protocol_classes} summarises the key characteristics of the three protocol
classes.

\begin{table}[ht]
\centering
\caption{Comparison of routing protocol classes for WMNs.}
\label{tab:protocol_classes}
\begin{tabular}{@{}lccc@{}}
\toprule
\textbf{Property} & \textbf{Proactive} & \textbf{Reactive} & \textbf{Hybrid} \\
\midrule
Route availability    & Always available & Discovery delay   & Zone: immediate; outside: delayed \\
Control overhead      & High (periodic)  & Low (on-demand)   & Medium \\
Memory usage          & High (full table)& Low (active only) & Medium \\
Scalability           & Poor in dense    & Better            & Good with tuning \\
Delay predictability  & High             & Low (first packet)& Medium \\
Energy efficiency     & Lower            & Higher            & Medium \\
Best for              & Small, latency-critical & Large, infrequent traffic & Large, mixed traffic \\
\bottomrule
\end{tabular}
\end{table}

\section{DSDV: Destination-Sequenced Distance Vector}

\subsection{Overview and History}

DSDV was proposed by Perkins and Bhagwat in 1994~\cite{perkins1994dsdv} as an
adaptation of the classic Bellman-Ford distance-vector algorithm for mobile ad~hoc
networks. The key innovation was the introduction of sequence numbers associated with
each route entry: every destination maintains a monotonically increasing sequence number,
and route updates bearing higher sequence numbers supersede older entries. This prevents
the \emph{counting to infinity} problem that plagues standard Bellman-Ford in dynamic
topologies.

\subsection{Routing Table Structure}

Every DSDV node maintains a routing table with one entry per known destination. Each
entry contains:

\begin{itemize}
  \item \textbf{Destination address:} the IP or MAC address of the destination node.
  \item \textbf{Next hop:} the address of the neighbouring node to which packets should be forwarded.
  \item \textbf{Metric:} typically the hop count to the destination, though other metrics (delay, ETX) can be used.
  \item \textbf{Sequence number:} an even number assigned by the destination when it
    initiates an update, or an odd number assigned by an intermediate node to mark a
    broken route. Higher is more recent; odd numbers signal link failure.
  \item \textbf{Install time:} the time at which the entry was last updated, used to expire stale routes.
\end{itemize}

\subsection{Update Mechanism}

DSDV uses two types of update messages:

\paragraph{Full dump.}
A complete routing table broadcast, sent periodically (typically every 30--60~seconds)
or when the routing table changes substantially. Full dumps allow newly joining nodes to
acquire the complete network topology quickly.

\paragraph{Incremental update.}
A partial table broadcast containing only the entries that have changed since the last
full dump. Incremental updates are sent more frequently (typically every 15~seconds) to
propagate link-state changes rapidly. Their smaller size reduces channel load between
full dumps.

Upon receiving an update, a node accepts a new entry if:
\begin{enumerate}
  \item The incoming sequence number is greater than the stored one (newer information), regardless of metric; or
  \item The incoming sequence number equals the stored one and the incoming metric is better.
\end{enumerate}

\subsection{Convergence and Stale Routes}

After a link failure, the node that detects the break immediately invalidates all routes
through that link (setting their sequence numbers to odd values indicating breakage) and
broadcasts an update. However, other nodes in the network may continue to use the broken
route for up to one full update cycle before receiving the invalidation. During this
convergence window, packets can be sent into a black hole.

In the INET~4.5 implementation used in this thesis, DSDV parameters are set to their
default values (full dump interval 30~s, incremental update interval 1~s). At $N = 10$
and a $1000\,\text{m} \times 1000\,\text{m}$ area, many source-destination pairs involve
multiple hops, and the convergence lag is long enough to cause significant packet drops,
as demonstrated in Section~\ref{sec:drop_analysis}.

\subsection{DSDV Algorithm Summary}

\begin{algorithm}[ht]
\caption{DSDV Route Update Processing}
\begin{algorithmic}[1]
\For{each entry $e$ in received update message}
  \If{$e.\text{seqNum} > \text{table}[e.\text{dest}].\text{seqNum}$}
    \State Accept: update table entry with new next-hop, metric, seqNum
  \ElsIf{$e.\text{seqNum} = \text{table}[e.\text{dest}].\text{seqNum}$}
    \If{$e.\text{metric} < \text{table}[e.\text{dest}].\text{metric}$}
      \State Accept: update table entry with better metric
    \Else
      \State Reject: same sequence, worse or equal metric
    \EndIf
  \Else
    \State Reject: older information (lower sequence number)
  \EndIf
\EndFor
\end{algorithmic}
\end{algorithm}

\section{AODV: Ad hoc On-Demand Distance Vector}

\subsection{Overview and Standardisation}

AODV was developed by Perkins, Belding-Royer, and Das~\cite{perkins2003aodv} and was
published as IETF RFC~3561 in July~2003. It extends DSDV's sequence-number mechanism
to a reactive framework, maintaining only the routes that are currently in active use.
AODV is one of the most widely implemented routing protocols in MANET and WMN research
and has been deployed in commercial products including Batman-Advanced and various IoT
gateway platforms.

\subsection{Route Discovery}

When a source node $S$ needs to send data to destination $D$ and has no valid routing
table entry for $D$, it broadcasts a Route Request (RREQ) message. The RREQ contains:

\begin{itemize}
  \item Source IP address and sequence number.
  \item Destination IP address and the last-known destination sequence number.
  \item Broadcast ID (incremented for each new RREQ from the source).
  \item Hop count (initialised to 0, incremented by each forwarder).
  \item Time-to-Live (TTL) field to limit flood scope.
\end{itemize}

Each intermediate node $M$ that receives the RREQ checks whether it has a valid,
sufficiently fresh route to $D$ (i.e., one with a sequence number at least as large as
the one in the RREQ). If so, $M$ generates a Route Reply (RREP) and sends it unicast
back to $S$ along the reverse path. Otherwise, $M$ increments the hop count, records
the reverse-path entry for $S$, and re-broadcasts the RREQ. Duplicate RREQ messages
(identified by source address + broadcast ID) are discarded.

The RREP travels back to $S$ along the reverse path established during the RREQ flood.
Each intermediate node that forwards the RREP creates a forward route entry for $D$
with the sequence number and hop count from the RREP. Once $S$ receives the RREP,
it begins data transmission immediately. The entire process, from RREQ to first data
packet, typically takes 2--4 round-trip times in a well-connected network.

\subsection{Route Maintenance}

Once a route is established, AODV uses two mechanisms to detect and respond to link
failures:

\paragraph{HELLO messages.}
Nodes periodically broadcast short HELLO messages to their neighbours. If no HELLO is
received from an active next-hop for a configurable interval, the link is assumed to have
failed.

\paragraph{Route Error (RERR).}
When a node detects that a link to an active next-hop has failed (either through missing
HELLOs or through a layer-2 notification from the MAC), it generates a Route Error
(RERR) message listing all destinations that become unreachable as a result. The RERR
is propagated upstream toward all sources using the broken route, which in turn may
initiate new route discoveries.

\subsection{Sequence Number Management}

AODV uses destination sequence numbers to ensure loop-free routes and to allow
comparisons between competing route options. Each destination node increments its own
sequence number before initiating a route reply. A route with a higher destination
sequence number is always preferred over one with a lower sequence number, regardless
of hop count.

\subsection{AODV Algorithm Summary}

\begin{algorithm}[ht]
\caption{AODV Route Discovery at Source Node}
\begin{algorithmic}[1]
\Procedure{SendData}{dest $D$, packet $P$}
  \If{valid route to $D$ exists in table}
    \State Forward $P$ to next-hop for $D$
  \Else
    \State Buffer $P$ in data queue
    \State Increment own RREQ sequence counter
    \State Broadcast $\text{RREQ}(\text{src}=S,\; \text{dest}=D,\; \text{hopCount}=0,\; \text{TTL}=\text{NET\_DIAMETER})$
    \State Start route-discovery timer
    \If{RREP received within timeout}
      \State Install route from RREP
      \State Flush buffered packets to $D$
    \Else
      \State Retry with expanded TTL (expanding ring search)
      \If{still no RREP}
        \State Drop buffered packets, notify application
      \EndIf
    \EndIf
  \EndIf
\EndProcedure
\end{algorithmic}
\end{algorithm}

\section{IEEE~802.11 DCF and MAC Layer Contention}

\subsection{Distributed Coordination Function (DCF)}

The IEEE~802.11 Distributed Coordination Function (DCF) is the fundamental access
control mechanism for WLANs and is the MAC protocol used in all nodes in this
simulation. DCF implements CSMA/CA: before transmitting, a node senses the channel.
If the channel is idle for a DCF Inter-Frame Space (DIFS) period, the node may
transmit. If the channel is busy, the node enters a random back-off period to reduce
the probability of collision.

\subsection{Binary Exponential Back-off}

The back-off duration is chosen uniformly at random from $[0, \text{CW})$ slots, where
CW is the Contention Window. The back-off timer counts down only while the channel is
idle; it is frozen when the channel becomes busy. After a collision (indicated by the
absence of an ACK from the receiver within a SIFS + ACK timeout), the sender doubles
the CW value:

\begin{equation}
  \text{CW}_{k+1} = \min\!\left(2 \cdot \text{CW}_k,\; \text{CW}_{\max}\right)
  \label{eq:beb}
\end{equation}

This is the Binary Exponential Back-off (BEB) algorithm. For IEEE~802.11g in mixed
mode:

\begin{align*}
  \text{CW}_{\min} &= 31 \text{ slots} \\
  \text{CW}_{\max} &= 1023 \text{ slots} \\
  \text{SIFS}       &= 10\;\mu\text{s} \\
  \text{DIFS}       &= 28\;\mu\text{s} \\
  \text{Slot time}  &= 9\;\mu\text{s}
\end{align*}

\subsection{Retry Limit and Frame Dropping}

IEEE~802.11 imposes a maximum retry count on every frame transmission. For the
802.11g standard (as configured in this simulation), the retry limit is 7~attempts. If
the frame has not been acknowledged after 7~transmission attempts, the MAC layer drops
the frame and generates a transmission failure notification to the upper layer. This is
the mechanism responsible for the packet drops observed in DSDV at $N = 50$, as
discussed in Section~\ref{sec:mac_evidence}.

The probability of a frame reaching the retry limit is negligible at low load, but
increases sharply as the offered load approaches the channel's saturation point.
Bianchi's seminal model~\cite{bianchi2000} characterises the saturation throughput of
802.11 DCF as:

\begin{equation}
  S = \frac{P_{tr} \cdot P_s \cdot E[L]}
         {(1 - P_{tr})\sigma + P_{tr}P_s T_s + P_{tr}(1-P_s)T_c}
  \label{eq:bianchi}
\end{equation}

where $P_{tr}$ is the probability that at least one node transmits in a slot, $P_s$ is
the probability of a successful transmission given that at least one node transmits,
$E[L]$ is the average payload size, $\sigma$ is the slot time, and $T_s$, $T_c$ are
the average times consumed by successful and collision events, respectively.

As the number of competing nodes $N$ increases, $P_{tr} \to 1$ and collisions dominate.
This is why DSDV's periodic routing broadcasts --- which add additional contending nodes
to the channel --- cause disproportionate performance degradation at $N = 50$ compared
with AODV's on-demand traffic model.

\section{Related Work}

A substantial body of literature compares AODV and DSDV under various conditions.
Das~et~al.~\cite{das2000performance} compared AODV and DSDV using the GloMoSim
simulator and found that AODV outperforms DSDV in terms of packet delivery ratio under
mobility, while DSDV achieves lower delay in static scenarios. Their work used
random-waypoint mobility and did not analyse the IEEE~802.11 MAC-layer behaviour.

Shukla and Mahajan~\cite{shukla2013performance} evaluated AODV, DSDV, and DSR in NS-2
and reported that AODV provides the best packet delivery ratio in medium-density networks,
while DSDV's routing overhead becomes problematic at higher densities. However, their
traffic model used a random source-destination pairing rather than the many-to-one IoT
pattern studied here.

Boukerche~\cite{boukerche2004} studied MANET routing protocols under varying network
sizes and traffic loads, finding a performance crossover between proactive and reactive
protocols at a density that depends on the traffic load. He did not provide MAC-layer
evidence for the observed crossover.

Bianchi's model~\cite{bianchi2000} provides the theoretical foundation for understanding
how the number of competing nodes affects IEEE~802.11 saturation throughput, and is used
in Appendix~\ref{app:math} to contextualise the MAC-level observations in this thesis.

Chatzimisios~et~al.~\cite{chatzimisios2003} studied the effect of the retry limit on
IEEE~802.11 throughput under high load conditions, confirming that the 7-retry abort
mechanism causes throughput collapse when the offered load exceeds the channel's
saturation point.

The contribution of this thesis relative to prior work is: (1) the use of a high-fidelity
IEEE~802.11 MAC model in OMNeT++/INET~4.5 rather than simplified MAC models in older
simulators; (2) the many-to-one traffic pattern representative of IoT data collection;
and (3) the explicit extraction of MAC-layer contention statistics to provide a causal
explanation of the observed performance differences.
